Para comprobar la puesta en marcha de Netdata en la Raspberry Pi se deben
ejecutar las siguientes pruebas desde una terminal. Cada prueba incluye el
resultado esperado que permite considerarla satisfactoria.

\section{Pruebas de puesta en marcha}

\subsection{Estado del servicio}

Comprobar que el servicio está habilitado y en ejecución:

\begin{lstlisting}[style=terminal]
sudo systemctl is-enabled netdata
sudo systemctl is-active netdata
sudo systemctl status netdata --no-pager
\end{lstlisting}

\textbf{Resultado esperado:} los dos primeros comandos muestran
\texttt{enabled} y \texttt{active}; el estado del servicio indica
\texttt{Active: active (running)}.

\subsection{Escucha del puerto de Netdata}

Verificar que Netdata escucha en el puerto TCP 19999:

\begin{lstlisting}[style=terminal]
sudo ss -lntp | grep ':19999'
\end{lstlisting}

\textbf{Resultado esperado:} aparece una línea en estado \texttt{LISTEN}
asociada al puerto \texttt{19999}.

\subsection{Respuesta de la interfaz web}

Consultar la interfaz desde la propia Raspberry Pi:

\begin{lstlisting}[style=terminal]
curl -sS -o /dev/null -w '%{http_code}\n' http://127.0.0.1:19999/
\end{lstlisting}

\textbf{Resultado esperado:} el código HTTP obtenido es \texttt{200}.

\subsection{Respuesta de la API}

Comprobar que la API entrega información del agente:

\begin{lstlisting}[style=terminal]
curl -fsS http://127.0.0.1:19999/api/v1/info | head
\end{lstlisting}

\textbf{Resultado esperado:} se recibe una respuesta JSON con información de
Netdata y del nodo, sin errores de conexión.

\subsection{Recolección de métricas}

Solicitar muestras recientes de uso de CPU mediante la API:

\begin{lstlisting}[style=terminal]
curl -fsS 'http://127.0.0.1:19999/api/v1/data?chart=system.cpu&after=-60&points=5'
\end{lstlisting}

\textbf{Resultado esperado:} se recibe una respuesta JSON con marcas de tiempo
y valores numéricos de CPU. Esto confirma que el agente no solo responde, sino
que también recolecta métricas.

\subsection{Registro de errores del servicio}

Revisar mensajes recientes del servicio durante el arranque:

\begin{lstlisting}[style=terminal]
sudo journalctl -u netdata -b --priority=err --no-pager
\end{lstlisting}

\textbf{Resultado esperado:} no aparecen errores críticos que impidan la
recolección o publicación de métricas.

\subsection{Arranque automático}

Reiniciar únicamente el servicio y repetir las comprobaciones básicas:

\begin{lstlisting}[style=terminal]
sudo systemctl restart netdata
sudo systemctl is-active netdata
curl -sS -o /dev/null -w '%{http_code}\n' http://127.0.0.1:19999/
\end{lstlisting}

\textbf{Resultado esperado:} el servicio vuelve a estado \texttt{active} y la
interfaz responde con código HTTP \texttt{200}. Para validar el arranque tras
un reinicio completo de la Raspberry Pi, reiniciar el equipo y repetir estas
dos últimas comprobaciones.

\section{Criterio de aceptación}

La puesta en marcha se considera satisfactoria cuando el servicio queda
habilitado y activo, escucha en el puerto 19999, responde por HTTP y por la API,
entrega métricas recientes y no registra errores críticos de operación.
